\documentclass[preview, border=4pt]{standalone}
\usepackage[a4paper, total={6in, 10in}]{geometry}
\usepackage{blindtext}
\usepackage{amssymb}
\usepackage{mathtools}
\begin{document}
\boldmath{Zeige $P = NP$}

Finde ein Problem P, die eindeutige Schneeflocke.

Gegeben: Eine Schneeflocke mit 5 gleichlangen Armen. \newline
Frage: Wie wird diese Schneeflocke eindeutig unter allen bereits vorhandenen Schneeflocken? \newline
Idee: Finde einen Algorithmus, der in polynomieller Zeit immer eine neue Schneeflocke erzeugt.

IB: Sei $S$ die Menge aller eindeutigen Schneeflocken, oBdA.\newline
IA: $|S| = 0$ trivial, $|S| = 1$ trivial\newline
IS: $|S| \rightarrow |S| + 1$, $|S| = n, n \in \mathbb{N}$\newline
Ziel: Mache aus 8 Knoten wieder 6 Knoten mit gleicher Kantenlänge in polynomieller Zeit.\newline
Idee: Berechne dazu die Restklasse 6.\newline
Jeder Knoten berechnet eine Funktion $f(b_1, \ldots, b_n) = \{0, \ldots, 9\}$, $b_i \in \{0, 1\}$.\newline
Bei Hinzufügen zwei neuer Knoten erhöht sich die Wertemenge des vorhandenen Graphen um maximal $18\equiv 0 \mod 6$.

Man dachte, das Problem P sei NP-vollständig. Da dieses Problem aber in polynomieller Zeit lösbar ist, ist es echt in $P$. 
\end{document}